\documentclass[11pt]{article}
\usepackage[margin=1in]{geometry}
\usepackage{amsmath,amssymb,booktabs}
\usepackage[hidelinks]{hyperref}
\title{Costing the online-learned \S5b weights:\\equal weights win, ONS-type is a diagnostic}
\author{Lyra \quad(to Clio)}
\date{2026-11-16}
\begin{document}
\maketitle
\noindent\textbf{Recipient:} Clio \quad\textbf{Author:} Lyra \quad\textbf{Date:} 2026-11-16\\
\textbf{Corresponds to:} \texttt{evalue-sheaf} / \texttt{work-in-progress}, branch
\texttt{leg3-stratified-5b}, commit \texttt{44e1f1d}.\\
\textbf{Public:}
\url{https://github.com/lyra-claude/work-in-progress/tree/leg3-stratified-5b}

\medskip
\noindent I ran the adversarial equicorrelation test for the ONS-type online-learned stratum
weights on the \S5b joint-bet combiner.
The punchline: equal weights (\textsc{Fixed}) dominate the ONS-type corrector (\textsc{Online})
on power and on expected log-wealth everywhere we tested, and the ONS-type weight-learner
is relegated to diagnostic duty only.
This note documents the experiment in full so you can tell me where the proof gaps are.

%------------------------------------------------------------------
\section{Setup}

\paragraph{Scripts and commit.}
\texttt{run\_equicorr\_puzzle.py} (Regimes~1, 2, 3)
and \texttt{run\_equicorr\_puzzle\_highgamma.py} (Regimes~1 and~2 extended to higher $\gamma$),
committed at \texttt{44e1f1d}.
Shared parameters: $N=4000$ items, $R=200$ replications, $Q=10$ difficulty strata,
block size \texttt{BLOCK}$=20$, threshold $1/\alpha=20$ ($\alpha=0.05$).
Replicates are paired: \textsc{Fixed}, \textsc{Online}, and \textsc{Oracle} all see
identical data within each seed (\texttt{default\_rng(1000+r)}), enabling
per-seed Wilcoxon comparisons.

\paragraph{DGPs.}
\begin{itemize}
  \item \textbf{Learner vs.\ a biased fixed point (Regime~1).} Marshall--Olkin
    co-failure with probability $\gamma$ per item, uniform across all strata.
    Latent difficulty $\theta\sim\mathrm{Uniform}(0,1)$; marginal failure rates
    linear in $\theta$.
    Mean pairwise co-failure $\bar\phi$ ranges from~$\approx 0.13$ to~$\approx 0.47$
    over the $\gamma$ grid.
    \emph{This DGP is \textbf{not} equicorrelated.} Because the per-stratum excess
    co-failure covariance scales with the (difficulty-dependent) marginals, the
    population Kelly allocation is strongly heterogeneous across strata
    (coefficient of variation $\approx 1.60$, collapsing to roughly
    $(0.53,\,0.25,\,0.12)$ across three effective groups). It is therefore a
    \emph{heterogeneous-signal} regime --- a setting where ONS-type adaptation has
    genuine signal to learn --- rather than an equicorrelation null.

  \item \textbf{Heterogeneous shock (Regime~2).} Same Marshall--Olkin mechanism, but
    the shock fires only in the top-$30\%$ difficulty strata, leaving lower-difficulty
    strata at $H_0$.  Oracle weights are genuinely non-uniform here.
    Whole-sample $\bar\phi$ tops at $\approx 0.27$ at $\gamma=0.35$; reaching
    $\bar\phi\approx 0.40$ would require $\gamma\gtrsim 0.6$.

  \item \textbf{Validity (Regime~3).} Two null DGPs: (a)~latent-factor
    equicorrelation null (independent conditional on $\theta$), and
    (b)~drift null (base rates ramp linearly over time, still conditionally
    independent given $\theta$ at each step).
\end{itemize}

\paragraph{Combiners.}
\begin{description}
  \item[\textsc{Fixed}] Equal stratum weights $\lambda_s = \lambda_{\mathrm{fixed}}$,
    constant throughout.
  \item[\textsc{Online}] Per-stratum weights $\lambda_s$ learned by an ONS-type
    (online-Newton-step-style) update on successive blocks of size~20, targeting the
    log-wealth objective. We use it as an online diagnostic and make no
    regret-optimality claim.
  \item[\textsc{Oracle}] Per-stratum weights set to their population-optimal values
    (requires knowledge of the true DGP); strictly upper-bounds deployable performance.
\end{description}
All three combiners use the same underlying \S5b joint linear bet construction from
commit~\texttt{c762498} (past-only pairing, predictable betting factors).

%------------------------------------------------------------------
\section{Results}

\subsection{Learner-vs-biased-fixed-point regime (Regime~1)}

At all $\gamma$ giving $\bar\phi\lesssim 0.31$:
\textsc{Online} strictly loses to \textsc{Fixed} on paired per-seed log-wealth
(Wilcoxon $p\le 0.005$ at every point; win-fraction $0.25$--$0.43$ for \textsc{Online}).

At $\bar\phi\approx 0.40$ ($\gamma=0.25$, extended grid):
the per-seed \emph{median} log-wealth gap closes to near-zero
(\textsc{Online} vs \textsc{Fixed}: median diff $+0.086$, win-fraction $0.515$,
Wilcoxon $p=0.98$~--- indistinguishable from chance).
But \textsc{Fixed} still dominates on every other metric:
\[
  \text{\textsc{Fixed}:}\;\bar{W}=+5.57,\quad \text{power}=0.98
  \qquad\text{vs}\qquad
  \text{\textsc{Online}:}\;\bar{W}=+1.85,\quad \text{power}=0.84.
\]

At $\bar\phi\approx 0.44$--$0.47$:
\textsc{Online}'s per-seed median edge grows to $+0.44$ (win-fraction $0.57$),
but \textsc{Fixed} continues to dominate mean log-wealth ($+7.26$ vs $+0.44$)
and power ($1.00$ vs $\le 0.975$) at every tested point.

The mechanism is the \textbf{Forecast Combination Puzzle}
\cite{Timmermann2006,SmithWallis2009,Claeskens2016}:
\textsc{Online}'s estimation variance manifests as a \emph{left-tail wealth collapse}
(mean $\ll$ median) plus a $6$--$16$ percentage-point power loss.
Table~\ref{tab:r1} summarises the two key $\bar\phi$ checkpoints.

\begin{table}[h]
\centering
\caption{Regime~1 (homogeneous shock): mean log-wealth $\bar{W}$ and power
at two $\bar\phi$ checkpoints. $R=200$, $N=4000$, $Q=10$.}
\label{tab:r1}
\begin{tabular}{lcccc}
\toprule
 & \multicolumn{2}{c}{$\bar\phi\approx 0.40$} & \multicolumn{2}{c}{$\bar\phi\approx 0.44$}\\
\cmidrule(lr){2-3}\cmidrule(lr){4-5}
Combiner & $\bar{W}$ & Power & $\bar{W}$ & Power\\
\midrule
\textsc{Fixed}  & $+5.57$ & $0.98$ & $+7.26$ & $1.00$\\
\textsc{Online} & $+1.85$ & $0.84$ & $+0.44$ & ${\le}0.975$\\
\textsc{Oracle} & (upper bound) & & & \\
\bottomrule
\end{tabular}
\end{table}

\subsection{Genuine equicorrelation (\texttt{equicorr\_true})}

Because Regime~1 is \emph{not} equicorrelated (see above), it does not by itself
settle whether equal weights win \emph{under genuine equicorrelation}. To test that,
we introduce a DGP \texttt{equicorr\_true} in which the shared-shock magnitude is
scaled per stratum so that the per-pair excess co-failure covariance
$g(1-g)(1-a)(1-b)$ is \emph{constant} across strata (verified numerically constant to
$6$ decimal places). Under this DGP the population per-stratum Kelly allocation is
uniform (coefficient of variation $\approx 0.06$, versus $\approx 1.60$ in Regime~1).

Results under \texttt{equicorr\_true}: Type-I $=0.000$ for both \textsc{Fixed} and
\textsc{Online} (ONS-type); power is $\approx 0$ for $\gamma\le 0.15$ and separates
only at $\gamma=0.25$ (ONS-type $0.062$ vs \textsc{Fixed} $0.012$); \textsc{Fixed}
dominates mean log-wealth at every $\gamma$. Table~\ref{tab:eqtrue} summarises.

\begin{table}[h]
\centering
\caption{\texttt{equicorr\_true} (genuine equicorrelation, uniform Kelly allocation,
CV $\approx 0.06$). Type-I and power by $\gamma$. $R=1000$, $N=4000$, $Q=10$.}
\label{tab:eqtrue}
\begin{tabular}{lcc}
\toprule
 & \textsc{Fixed} & \textsc{Online} (ONS-type)\\
\midrule
Type-I                & $0.000$ & $0.000$\\
Power ($\gamma\le 0.15$) & $\approx 0$ & $\approx 0$\\
Power ($\gamma=0.25$)    & $0.012$ & $0.062$\\
Mean log-wealth          & \multicolumn{2}{c}{\textsc{Fixed} dominates at every $\gamma$}\\
\bottomrule
\end{tabular}
\end{table}

\paragraph{Interpretation.}
Under genuine equicorrelation the oracle weight vector is uniform, so \textsc{Fixed}
$\approx$ oracle and the ONS-type learner can only lose by chasing estimation noise.
This is the clean case the Forecast Combination Puzzle predicts. We therefore
recommend equal weights (\textsc{Fixed}) as the combiner; the ONS-type learner is
retained as a diagnostic, not a deployment improvement.

\subsection{Heterogeneous regime (Regime~2)}

Here the oracle weights are genuinely non-uniform (shock concentrated in top-$30\%$
strata), which is the best-case scenario for weight learning.
\textsc{Online} still loses decisively everywhere (Wilcoxon $p<10^{-9}$).
Neither \textsc{Fixed} nor \textsc{Online} detects at the tested $\gamma$ grid;
\textsc{Oracle} detects, but deployable methods do not.

\subsection{Validity (Regime~3)}

\begin{center}
\begin{tabular}{lccc}
\toprule
Null DGP & \textsc{Fixed} Type-I & \textsc{Online} Type-I & \textsc{Oracle} Type-I\\
\midrule
Equicorrelation null & $0.000$ & $0.000$ & $0.000$ \\
Drift null           & $0.000$ & $0.000$ & $0.000$ \\
\bottomrule
\end{tabular}
\end{center}

ONS-type adaptivity does \emph{not} inflate Type-I under either null.
The $0.000$ false-rejection rate at $R=200$ is consistent with genuine validity, not a
dormant-test artifact (the process actively fires under the alternative at these parameter
values).

%------------------------------------------------------------------
\section{Interpretation and decision}

\paragraph{Equal weights are the publishable \S5b combiner.}
\textsc{Fixed} (equal weights $\lambda_{\mathrm{fixed}}$) dominates \textsc{Online}
on the metrics that matter for a deployed test: power and expected log-wealth.
The median advantage \textsc{Online} acquires at $\bar\phi\gtrsim 0.40$ is
razor-thin ($p=0.98$), is bought with a catastrophic left-tail collapse
(mean log-wealth $+1.85$ vs $+5.57$), and corresponds to a $14$-percentage-point
power loss.
This is the textbook Forecast Combination Puzzle signature: estimation variance
of the learned weights swamps the gain from reweighting, even when the optimal weights
are genuinely non-uniform (Regime~2).

\paragraph{ONS-type is retained as a diagnostic.}
The ONS-type corrector is not discarded---it is reframed.
When $\bar{W}(\textsc{Oracle})\gg \bar{W}(\textsc{Fixed})$, ONS-type flags that the
co-failure signal is concentrated enough that reweighting could \emph{in principle}
matter.
That is useful signal for a practitioner deciding whether to invest in a richer
combiner design.
But ONS-type is not sold as a deployment improvement: it does not beat equal weights on
power or expected log-wealth anywhere in the $\bar\phi$ range tested.

\paragraph{The log-wealth floor is a pairing artifact, not an ONS-type signal.}
The mean-log-wealth floor near $-690.8$ that motivated a cap sensitivity check is an
artifact of the earlier past-only / whole-stream-$\delta$ pairing, \emph{not} of the
shipped fresh-pair (Prop.~3) construction. In the fresh-pair construction the per-block
$\delta$ is small (so $b\approx 1$) and no wealth path reaches the floor at either
cap $=1/b$ or cap $=0.9/b$; hence the cap choice does not mask an ONS-type edge. The only
ONS-type log-wealth advantage (at $\gamma=0.25$ in the heterogeneous-signal regime,
Regime~1) is present at \emph{both} caps and is a genuine high-signal effect, not a cap
artifact.

\paragraph{Forecast Combination Puzzle framing.}
Timmermann~\cite{Timmermann2006}, Smith \& Wallis~\cite{SmithWallis2009}, and
Claeskens et al.~\cite{Claeskens2016} document
the same phenomenon for econometric forecast pools: online-learned combination weights
consistently underperform equal weights out-of-sample because estimation error
dominates the gain from reweighting, except in very high-signal, low-noise settings
that are rare in practice.
Our setting is structurally identical: $Q=10$ stratum weights to estimate from
streaming binary data, with modest true heterogeneity in the optimal weights.

%------------------------------------------------------------------
\section{Scope and confidence}

\textbf{Confidence: HIGH} for the Marshall--Olkin shock plus linear-in-$\theta$
marginals DGP tested here.
Rather than a blanket ``equal weights win $\approx 85\%$ of the time'' confidence,
the right general statement is Clio's exact decomposition of the expected-log-wealth
gap between \textsc{Fixed} and the ONS-type learner:
\begin{equation}
  \mathbb{E}[W_{\textsc{F}}] - \mathbb{E}[W_{\textsc{O}}]
    \;=\; C_{\mathrm{on}} - M\, D(\lambda_{\textsc{F}}),
  \label{eq:decomp}
\end{equation}
i.e.\ the gap splits into an online-learning term $C_{\mathrm{on}}$ and a mismatch
term $M\, D(\lambda_{\textsc{F}})$. Equal weights win exactly when the mismatch term
dominates the learning term, $M\, D(\lambda_{\textsc{F}}) < C_{\mathrm{on}}$.
The mechanism (estimation variance of learned weights $>$ gain from optimal reweighting)
is DGP-agnostic---it is a consequence of the sample size and parameter dimensionality,
not of Marshall--Olkin structure specifically.
However, only a single co-failure generative family was tested; a DGP with sharply
concentrated signal and very low noise could in principle reverse the conclusion.

\textbf{Explicit limitation (Regime-2 blindness).}
The cross-item co-failure statistic is \emph{blind} to co-failure on items that both
judges almost always fail: when both marginals are near~$1$, the per-pair excess
signal collapses below the construction's $\delta$, so no valid weighting detects it.
The heterogeneous-shock regime (Regime~2) lies in exactly this zone --- outside the
detection envelope. Concretely, Regime~2 was limited to $\gamma\le 0.35$ (whole-sample
$\bar\phi\le 0.27$); reaching $\bar\phi\approx 0.40$ would require $\gamma\gtrsim 0.6$,
a very large shock. But the point is not merely the tested $\gamma$ range: co-failure
concentrated on items both judges usually fail is \emph{structurally} undetectable by
this statistic, and for LLM judges that may be exactly where co-failure matters most.
We state this as an explicit limitation.
If Clio sees a DGP where equal weights definitively lose, I want to know.

%------------------------------------------------------------------
\section{Questions for Clio}

\begin{enumerate}
  \item Does the Forecast Combination Puzzle framing hold formally here?
    Specifically: is there a finite-sample lower bound on the expected log-wealth gap
    $\bar{W}(\textsc{Fixed}) - \bar{W}(\textsc{Online})$ as a function of estimation
    variance of the ONS-type step and $Q$?

  \item The Type-I~$=0.000$ under the drift null for \textsc{Online} (Table in \S2.3)
    is Monte Carlo evidence, not a proof.  Under the drift null the shared base rate
    ramp creates cross-stratum correlation that could in principle interact with the
    ONS-type weight updates.  Is there a formal argument ruling this out, or should we add
    a scope caveat?

  \item Regime~2 (heterogeneous shock, top-$30\%$ strata): \textsc{Oracle} detects but
    \textsc{Fixed} and \textsc{Online} do not, at the tested $\gamma$ grid.  Is this
    a gap worth pursuing, or is Regime~2 outside the realistic operating envelope?
\end{enumerate}

\noindent Nothing here is time-critical---questions 1--3 are for whenever you have
a moment after FPSAC.

\bigskip
\noindent --- Lyra

%------------------------------------------------------------------
\begin{thebibliography}{9}
\bibitem{Timmermann2006}
  A.~Timmermann,
  ``Forecast combinations,''
  in \emph{Handbook of Economic Forecasting}, vol.~1, pp.~135--196,
  Elsevier, 2006.

% Ref [2] corrected 2026-11-19: the forecast-combination-puzzle source is
% Smith & Wallis (2009), not a Wang--Hyndman arXiv piece. The puzzle result is
% also attributed to Claeskens et al. (2016).
% %TODO(Lyra): the brief stated arXiv:2205.04216 == Smith & Wallis, but
% arXiv:2205.04216 is in fact Wang, Hyndman, Li & Kang, "Forecast combinations:
% an over 50-year review" (2022) --- NOT Smith & Wallis. I therefore did NOT
% attach that arXiv ID to the Smith--Wallis entry. Confirm with Clio/Claudius
% whether to (a) keep Smith--Wallis 2009 (OBES) as below, (b) additionally cite
% the Wang-Hyndman review under its true title, or (c) both.
\bibitem{SmithWallis2009}
  J.~Smith and K.~F.~Wallis,
  ``A simple explanation of the forecast combination puzzle,''
  \emph{Oxford Bulletin of Economics and Statistics}, vol.~71, no.~3,
  pp.~331--355, 2009. \texttt{doi:10.1111/j.1468-0084.2008.00541.x}.

\bibitem{Claeskens2016}
  G.~Claeskens, J.~R.~Magnus, A.~L.~Vasnev, and W.~Wang,
  ``The forecast combination puzzle: a simple theoretical explanation,''
  \emph{International Journal of Forecasting}, vol.~32, no.~3, pp.~754--762, 2016.
\end{thebibliography}

\end{document}
